\section{Motivation: recovery and decision-rule validation}
\label{sec:intro}

This paper revisits a specific validation claim: a single-cell analysis
reported that an estimator recovered an imposed effect and treated this as
validation of a 50-cell reporting cutpoint. The original design required
interval inclusion and discrimination, which the recovery ratio did not
measure. Section~\ref{sec:calibration} reproduces the claim and criteria,
then tests them directly. The empirical contribution is the failed calibration
and its sensitivity to sampling and aggregation. The algebra used to explain
the original comparison is elementary; the trial examples are worked checks
of established estimator relationships.

Our motivating application separates two explanations for differentiation-marker
loss in colorectal cancer: fewer marker-producing cells (a \emph{compositional}
change), or lower output from the cells that remain (an \emph{intrinsic}
change). These explanations suggest different biological interpretations.
An empirical comparison of per-cell means requires the relevant population
to be observed in both tissues. Let $f_N,f_T$ denote the relevant cell fractions
in reference and diseased tissue, and $m_N,m_T$ their mean per-cell output.
A reference-based decomposition, related to demographic standardisation and
Oaxaca--Blinder decompositions \citep{kitagawa1955,oaxaca1973,blinder1973}, is
\begin{equation}
\underbrace{(f_T-f_N)m_N}_{\text{compositional}} +
\underbrace{f_N(m_T-m_N)}_{\text{intrinsic}} +
\underbrace{(f_T-f_N)(m_T-m_N)}_{\text{interaction}}
= f_Tm_T-f_Nm_N.
\label{eq:kitagawa}
\end{equation}
The interaction is reported separately. If no relevant cells are observed in
an arm, its sample mean is undefined; the identity cannot supply the missing
mean. In a synthetic population with no such cells, the corresponding
conditional mean itself has no referent. Substituting zero in either case
turns an unavailable contrast into an apparent effect.

The implemented reporting rule returns an estimate at $n\geq50$, flags a wide
interval at $20\leq n<50$, and returns \notest{} below 20, where $n$ is the
relevant-cell count in diseased tissue. Unavailable reported contrasts are
stored as \texttt{None}, not zero. The diagnostic oracle calculations examined
below also retain raw arithmetic before this reporting gate. The question is
whether the earlier recovery evidence established the interval-performance
requirements used to justify the cutpoints.

\paragraph{Relation to existing validation methods.}
Simulation-study methodology already separates estimands, methods and
performance measures \citep{burton2006,morris2019}. Generator--analysis alignment
is discussed in congeniality \citep{meng1994}, plasmode and semi-synthetic
benchmarks \citep{franklin2014,curth2021}, interventional evaluation
\citep{gentzel2019} and the inverse crime \citep{wirgin2004,kaipio2007}; generators
can also disadvantage good estimators \citep{shaw2026}. Our analysis applies these established distinctions to a documented
reporting-rule decision. For clinical trials, external-control validation
already examines bias and type-I error using held-out studies
\citep{ventz2019}. Prognostic adjustment with historical data retains trial
randomisation and evaluates precision gains \citep{schuler2022}; target-trial
emulation states the causal question and design explicitly \citep{hernan2016}.
These are different uses of a patient model and need different validation
experiments. Our case contributes an inspectable failure analysis relevant to
that design problem, not evidence of failure in those methods. Simulation-based
calibration evaluates posterior ranks \citep{talts2018,modrak2025}.
Appendix~\ref{app:relatedwork}
provides further context.
